\documentclass[12pt,a4paper]{article}

% ------------------------------
% Paquetes para idioma y tipografía
% ------------------------------
\usepackage[spanish,es-tabla]{babel}
\usepackage[utf8]{inputenc}
\usepackage[T1]{fontenc}
\usepackage{lmodern}
\usepackage{microtype}

% ------------------------------
% Paquetes para matemáticas y símbolos
% ------------------------------
\usepackage{amsmath}
\usepackage{amsfonts}
\usepackage{amssymb}

% ------------------------------
% Paquetes para formato y gráficos
% ------------------------------
\usepackage{geometry}
\geometry{top=2.5cm,bottom=2.5cm,left=2.5cm,right=2.5cm}
\usepackage{graphicx}
\usepackage{float}
\usepackage{array}
\usepackage{booktabs}
\usepackage{longtable}
\usepackage{setspace}
\usepackage{xcolor}
\usepackage{listings}
\usepackage{enumitem}

% ------------------------------
% Configuración de listings (SQL / PL-pgSQL / Bash)
% ------------------------------
\lstdefinelanguage{PostgreSQL}{
  keywords={SELECT,INSERT,UPDATE,DELETE,CREATE,ALTER,DROP,TABLE,VIEW,TRIGGER,
            TABLESPACE,LOCATION,INDEX,PARTITION,PARTITIONS,OF,BY,RANGE,VALUES,FROM,
            USER,ROLE,GRANT,REVOKE,TO,WHERE,AND,OR,IN,NOT,NULL,PRIMARY,KEY,
            FOREIGN,REFERENCES,CHECK,DEFAULT,ON,AS,FOR,EACH,ROW,STATEMENT,BEFORE,
            AFTER,BEGIN,END,DECLARE,EXCEPTION,WHEN,LOOP,WHILE,IF,THEN,ELSE,ELSIF,
            COMMIT,ROLLBACK,EXECUTE,EXPLAIN,ANALYZE,SET,SHOW,RETURNS,RETURN,
            FUNCTION,PROCEDURE,CALL,LANGUAGE,DO,RAISE,NOTICE,SCHEMA,POLICY,
            ENABLE,FORCE,LEVEL,SECURITY,DEFINER,USING,WITH,CHECK,TRUNCATE,
            REINDEX,VACUUM,ATTACH,DETACH,INCLUDE,ADD,JOIN,GROUP,ORDER,LIMIT,
            UNION,ALL,DISTINCT,CASE,END,ASC,DESC,FILTER,LOGIN,PASSWORD,
            CONNECT,DATABASE,USAGE,SEQUENCES,TABLES,IDENTITY,GENERATED,
            UNIQUE,LIKE,INCLUDING,DEFAULTS,CONSTRAINTS,SYSTEM,SCOPE,MINVALUE,
            MAXVALUE,NOLOGIN,AUTHORIZATION,EXTENSION,EXISTS,RESET,CONSTANT},
  sensitive=false,
  morecomment=[l]{--},
  morecomment=[s]{/*}{*/},
  morestring=[b]',
}

\lstset{
  language=PostgreSQL,
  basicstyle=\ttfamily\footnotesize,
  keywordstyle=\bfseries\color{blue!70!black},
  commentstyle=\itshape\color{gray!90!black},
  stringstyle=\color{orange!80!black},
  numbers=left,
  numberstyle=\tiny\color{gray},
  stepnumber=1,
  numbersep=8pt,
  breaklines=true,
  breakatwhitespace=true,
  frame=single,
  backgroundcolor=\color{gray!6},
  tabsize=2,
  captionpos=b,
  showstringspaces=false,
  extendedchars=true,
  literate=%
    {á}{{\'a}}1
    {é}{{\'e}}1
    {í}{{\'i}}1
    {ó}{{\'o}}1
    {ú}{{\'u}}1
    {Á}{{\'A}}1
    {É}{{\'E}}1
    {Í}{{\'I}}1
    {Ó}{{\'O}}1
    {Ú}{{\'U}}1
    {ñ}{{\~n}}1
    {Ñ}{{\~N}}1
    {¿}{{?`}}1
    {¡}{{!`}}1
    {~}{{$\sim$}}1
}

% Código SQL generado junto con cada evidencia (carpeta sql/)
\setlength{\emergencystretch}{3em}
\newcommand{\codigo}[2]{\lstinputlisting[caption={#2}]{sql/#1.sql}}

% Evidencia (captura de terminal generada con psql sobre el contenedor)
\newcommand{\evid}[3]{%
  \begin{figure}[H]
    \centering
    \includegraphics[width=\textwidth,height=0.86\textheight,keepaspectratio]{img/#1.png}
    \caption{#2}
    \label{fig:#3}
  \end{figure}}

% ------------------------------
% Paquetes para referencias y citas
% ------------------------------
\usepackage{csquotes}
\usepackage{hyperref}
\hypersetup{
  colorlinks=true,
  linkcolor=blue!70!black,
  urlcolor=blue!70!black,
  citecolor=blue!70!black,
}

% ==========================================
% PORTADA
% ==========================================
\begin{document}

\begin{titlepage}
    \centering

    \begin{minipage}{0.25\textwidth}
        \centering
        \includegraphics[width=3cm]{Escudo-UNAM.png}
    \end{minipage}
    \hfill
    \begin{minipage}{0.45\textwidth}
        \centering
        \textbf{\LARGE UNIVERSIDAD}\\[0.25cm]
        \textbf{\LARGE NACIONAL}\\[0.25cm]
        \textbf{\LARGE AUTÓNOMA}\\[0.25cm]
        \textbf{\LARGE DE MÉXICO}
    \end{minipage}
    \hfill
    \begin{minipage}{0.25\textwidth}
        \centering
        \includegraphics[width=3cm]{escudofi_negro.jpg}
    \end{minipage}

    \vspace{0.7cm}
    {\large\textbf{FACULTAD DE INGENIERÍA}}\\[1.2cm]

    {\Large\textbf{PROYECTO FINAL}}\\[0.4cm]
    {\LARGE\textbf{Diseño, Seguridad y Optimización de un}}\\[0.2cm]
    {\LARGE\textbf{Sistema de Gestión de Empleados}}\\[1.2cm]

    {\large\textbf{Materia:} Bases de Datos Avanzadas}\\[0.3cm]
    {\large\textbf{Grupo:} 3 \quad \textbf{Equipo:} 8}\\[1.2cm]

    \textbf{Integrantes:}\\[0.3cm]
    Rojas Chimal Karla Beatriz --- \texttt{320048025}\\[0.3cm]
    Calderón Gutiérrez Victor Emiliano --- \texttt{423097331}\\[1.2cm]

    \textbf{Profesor:} Mtro. Gerardo Gabriel Carrasco Zúñiga\\[0.5cm]

    \textbf{Fecha:} 22 de noviembre de 2026 \quad \textbf{Versión:} 1.0\\[0.8cm]

    \textbf{Calificación:} \_\allowbreak{}\_\allowbreak{}\_\allowbreak{}\_\allowbreak{}\_\allowbreak{}\_\allowbreak{}\_\allowbreak{}\_\allowbreak{}\_\allowbreak{}\_\allowbreak{}\_\allowbreak{}\_\allowbreak{}\_\allowbreak{}\_\allowbreak{}\_\allowbreak{}\_\allowbreak{}

    \vfill
\end{titlepage}

\spacing{1.3}

% ==========================================
% TABLA DE CONTENIDOS
% ==========================================
\tableofcontents
\newpage

% ==========================================
% 1. INTRODUCCIÓN
% ==========================================
\section{Introducción y objetivos}

\subsection{Objetivo}

Aplicar los conocimientos del semestre de Bases de Datos Avanzadas en un sistema de gestión de empleados que integre diseño con tablas particionadas y múltiples tablespaces, seguridad multinivel (RBAC, DAC y MAC), auditoría automática, optimización del rendimiento y administración del almacenamiento, documentando cada paso con evidencias reproducibles.

El enunciado está redactado para Oracle (SQL*Plus, \texttt{DBA\_\allowbreak{}ADVISOR\_\allowbreak{}*}, \texttt{DBMS\_\allowbreak{}XPLAN}). Igual que en las prácticas 1 a 4, el proyecto se implementa en \textbf{PostgreSQL 16 ejecutado en un contenedor Docker}; en la sección~\ref{sec:equivalencias} se indica cómo se tradujo cada elemento. Todas las cifras de este documento provienen de la ejecución real de los scripts (transcripción completa en \texttt{Dockerpfinal/transcripcion.txt}). Las evidencias tienen formato de terminal, pero no son capturas de pantalla: \texttt{generar\_\allowbreak{}evidencias.py} ejecuta los diez scripts sobre un contenedor recién creado y dibuja como imagen la salida real de \texttt{psql}, de modo que cualquiera puede regenerarlas y obtener los mismos costos y conteos (los milisegundos varían ligeramente con el equipo).

\subsection{Requerimientos y dónde se cubren}

\begin{table}[H]
\centering
\small
\renewcommand{\arraystretch}{1.3}
\begin{tabular}{p{0.52\textwidth}p{0.36\textwidth}}
\toprule
\textbf{Requerimiento del enunciado} & \textbf{Sección} \\
\midrule
Esquema de empleados, departamentos y roles con tablas particionadas y varios tablespaces (con su tamaño definido) & \ref{sec:diseno} \\
Usuarios admin/analista/auditor, roles, DAC, MAC y RBAC; pruebas de acceso, modificación y auditoría & \ref{sec:seguridad}, \ref{sec:pruebas} \\
Vista y trigger de auditoría y control de cambios & \ref{sec:auditoria} \\
Inserción masiva (mín. 50,000 registros) con bloque procedimental & \ref{sec:carga} \\
Análisis de 3 consultas complejas con \texttt{EXPLAIN} & \ref{sec:antes} \\
Recreación de índices y evaluación del costo & \ref{sec:indices} \\
Ajuste de parámetros (\texttt{DB\_\allowbreak{}BLOCK\_\allowbreak{}SIZE}, etc.) y advisor & \ref{sec:advisor}, \ref{sec:parametros} \\
Monitoreo de tablespaces, ajustes dinámicos y mantenimiento automático & \ref{sec:almacenamiento} \\
Comparativo antes/después y propuesta para producción & \ref{sec:comparativo}, \ref{sec:recomendaciones} \\
Guía de demostración con cambio de rol & \ref{sec:demo} \\
Conclusión de lo logrado y lo no cubierto & \ref{sec:conclusiones} \\
\bottomrule
\end{tabular}
\caption{Trazabilidad entre el enunciado y este documento.}
\end{table}

\subsection{Equivalencias Oracle $\rightarrow$ PostgreSQL}
\label{sec:equivalencias}

\begin{table}[H]
\centering
\small
\renewcommand{\arraystretch}{1.3}
\begin{tabular}{p{0.42\textwidth}p{0.47\textwidth}}
\toprule
\textbf{Oracle (enunciado)} & \textbf{PostgreSQL (adaptación)} \\
\midrule
VM VirtualBox, \texttt{sqlplus / as sysdba} & Contenedor Docker, \texttt{psql -U postgres} \\
\texttt{CREATE TABLESPACE ... DATAFILE ... AUTOEXTEND ON} & \texttt{CREATE TABLESPACE ... LOCATION} (sin datafiles ni \texttt{AUTOEXTEND}) \\
\texttt{CREATE USER ... IDENTIFIED BY} & \texttt{CREATE USER ... LOGIN PASSWORD} \\
\texttt{CREATE ROLE}, \texttt{GRANT} & \texttt{CREATE ROLE ... NOLOGIN}, \texttt{GRANT} \\
MAC (Oracle Label Security / VPD) & Niveles de autorización + \textit{Row Level Security} (\texttt{CREATE POLICY}) \\
\texttt{PARTITION BY RANGE ... VALUES LESS THAN} & \texttt{PARTITION BY RANGE} + \texttt{PARTITION OF ... FOR VALUES FROM ... TO} \\
Bloque PL/SQL \texttt{BEGIN ... END;} & Bloque \texttt{DO \$\$ ... \$\$} (PL/pgSQL) \\
\texttt{SYSDATE}, \texttt{SYS\_\allowbreak{}CONTEXT} & \texttt{CURRENT\_\allowbreak{}DATE}, \texttt{session\_\allowbreak{}user}, \texttt{inet\_\allowbreak{}client\_\allowbreak{}addr()} \\
\texttt{EXPLAIN PLAN FOR} + \texttt{DBMS\_\allowbreak{}XPLAN.DISPLAY} & \texttt{EXPLAIN} y \texttt{EXPLAIN (ANALYZE, FORMAT JSON)} \\
\texttt{DB\_\allowbreak{}BLOCK\_\allowbreak{}SIZE}, \texttt{SGA\_\allowbreak{}TARGET}, \texttt{PGA\_\allowbreak{}AGGREGATE\_\allowbreak{}TARGET} & \texttt{block\_\allowbreak{}size}, \texttt{shared\_\allowbreak{}buffers}, \texttt{work\_\allowbreak{}mem} \\
\texttt{ALTER SYSTEM SET ... SCOPE=BOTH} & \texttt{ALTER SYSTEM SET} + \texttt{pg\_\allowbreak{}reload\_\allowbreak{}conf()} \\
\texttt{DBA\_\allowbreak{}ADVISOR\_\allowbreak{}FINDINGS}, \texttt{DBA\_\allowbreak{}ADVISOR\_\allowbreak{}RECOMMENDATIONS} & Vistas propias sobre \texttt{pg\_\allowbreak{}stat\_\allowbreak{}user\_\allowbreak{}tables}, \texttt{pg\_\allowbreak{}stat\_\allowbreak{}user\_\allowbreak{}indexes} y \texttt{pg\_\allowbreak{}stat\_\allowbreak{}statements} \\
\texttt{DBA\_\allowbreak{}DATA\_\allowbreak{}FILES}, \texttt{DBA\_\allowbreak{}FREE\_\allowbreak{}SPACE} & \texttt{pg\_\allowbreak{}tablespace\_\allowbreak{}size()}, \texttt{pg\_\allowbreak{}tablespace\_\allowbreak{}location()} \\
\texttt{ALTER TABLE ... TRUNCATE PARTITION} & \texttt{DETACH PARTITION}, \texttt{TRUNCATE}, \texttt{ATTACH PARTITION} \\
\texttt{ALTER INDEX ... REBUILD} & \texttt{REINDEX INDEX} \\
\bottomrule
\end{tabular}
\caption{Equivalencia de comandos, parámetros y vistas del diccionario de datos.}
\end{table}

% ==========================================
% 2. ENTORNO
% ==========================================
\section{Entorno de trabajo con Docker}

El entorno se define de forma declarativa en \texttt{docker-compose\allowbreak{}.yaml} (carpeta \texttt{Dockerpfinal}), con la misma convención de las prácticas anteriores: imagen \texttt{postgres:16}, red y volúmenes con nombre. Dos particularidades respecto a la práctica 3:

\begin{itemize}
    \item El servicio auxiliar \texttt{init-tablespaces} (se ejecuta como \texttt{root} una sola vez) crea los cinco directorios de los tablespaces y los asigna al usuario \texttt{postgres}, por lo que ya no hay que preparar el volumen manualmente.
    \item El script \texttt{initdb/01-hba-scram.sh} deja las conexiones TCP de loopback con \texttt{scram-sha-256}. Así, las pruebas con \texttt{psql -h localhost} validan de verdad la contraseña de cada usuario, mientras que el socket Unix queda en \texttt{trust} (análogo a \texttt{/ as sysdba}).
\end{itemize}

Además se precarga \texttt{pg\_\allowbreak{}stat\_\allowbreak{}statements}, base del \textit{advisor} que se construye en la sección~\ref{sec:advisor}, y se monta la carpeta \texttt{sql/} con los diez scripts numerados del proyecto. El archivo siguiente muestra los dos servicios base; el tercero, \texttt{mantenimiento}, programa el mantenimiento automático y se describe en la sección~\ref{sec:mant_programado}.

\begin{lstlisting}[language={},caption={Archivo docker-compose.yaml del Proyecto Final}]
services:
  init-tablespaces:
    image: postgres:16
    container_name: bd-lab-final-init
    user: root
    command:
      - sh
      - -c
      - >
        mkdir -p /ts/data /ts/idx /ts/part /ts/audit /ts/hist &&
        chown -R postgres:postgres /ts && chmod 700 /ts/*
    volumes:
      - bd_tablespace_final:/ts
    restart: "no"

  postgres-final:
    image: postgres:16
    container_name: bd-lab-final
    command: ["postgres", "-c", "shared_preload_libraries=pg_stat_statements"]
    depends_on:
      init-tablespaces:
        condition: service_completed_successfully
    environment:
      POSTGRES_USER: postgres          # superusuario (equivalente a SYS / DBA)
      POSTGRES_PASSWORD: Admin12345
      POSTGRES_DB: rrhh                # BD del Sistema de Gestion de Empleados
      PGUSER: postgres
      PGDATABASE: rrhh
    ports:
      - "5432:5432"
    volumes:
      - bd_datos_final:/var/lib/postgresql/data
      - bd_tablespace_final:/var/lib/postgresql/tablespaces
      - ./sql:/sql:ro
      - ./initdb:/docker-entrypoint-initdb.d:ro
    networks:
      - red-final
    healthcheck:
      test: ["CMD-SHELL", "pg_isready -U postgres -d rrhh"]
      interval: 5s
      timeout: 3s
      retries: 10
    restart: unless-stopped

volumes:
  bd_datos_final:
    name: bd_datos_final
  bd_tablespace_final:
    name: bd_tablespace_final

networks:
  red-final:
    name: red-final
\end{lstlisting}

\begin{lstlisting}[language={},caption={Despliegue del entorno y conexión}]
docker compose up -d
docker ps
docker exec bd-lab-final pg_isready -U postgres -d rrhh
docker exec -it bd-lab-final psql          # PGUSER y PGDATABASE ya vienen definidos
\end{lstlisting}

\evid{ev01_entorno}{Despliegue del contenedor \texttt{bd-lab-final}, verificación con \texttt{pg\_\allowbreak{}isready} y versión del servidor.}{ev01}

% ==========================================
% 3. DISEÑO
% ==========================================
\section{Diseño avanzado de la base de datos}
\label{sec:diseno}

\subsection{Tablespaces}

Se crean cinco tablespaces, cada uno con un directorio propio dedicado a un tipo de carga: datos generales, índices, particiones activas, bitácora de auditoría y datos históricos. De esta forma cada tipo de objeto puede crecer, respaldarse o moverse a otro disco de forma independiente.

\codigo{ev02_tablespaces}{Creación de los cinco tablespaces del sistema}
\evid{ev02_tablespaces}{Tablespaces creados y listados con \texttt{\textbackslash db+}.}{ev02}

\paragraph{Dimensionamiento.} En Oracle el tamaño se fija en el \textit{datafile} (\texttt{SIZE}, \texttt{AUTOEXTEND ... MAXSIZE}); en PostgreSQL un tablespace es un directorio y no admite tamaño en \texttt{CREATE TABLESPACE}. Por eso el tamaño de cada uno se definió por cálculo, para un horizonte de \textbf{300,000 empleados} (cinco veces la carga actual), a partir de los bytes por fila \textbf{medidos} (4,452 + 3,020\,kB de particiones $\approx$ 6,972\,kB / 60,000 filas $\approx$ 120\,B por fila de datos y 6,772\,kB / 60,000 $\approx$ 115\,B por fila de índices). El tamaño asignado forma parte del \textbf{script de creación}: \texttt{01\_\allowbreak{}instancia\_\allowbreak{}roles.sql} lo registra en la tabla \texttt{capacidad\_\allowbreak{}tablespaces} justo después de crear el esquema (listado y evidencia~\ref{fig:ev04}), y la vista \texttt{v\_\allowbreak{}uso\_\allowbreak{}tablespaces} compara el uso real contra él y marca \texttt{ALERTA} al llegar al 80\,\% (sección~\ref{sec:almacenamiento}). PostgreSQL no impide escribir al pasar ese tamaño: en producción cada directorio se montaría en un volumen de ese tamaño, que es el límite físico equivalente a \texttt{MAXSIZE}.

\begin{table}[H]
\centering
\small
\renewcommand{\arraystretch}{1.3}
\setlength{\tabcolsep}{4pt}
\begin{tabular}{lp{0.24\textwidth}p{0.25\textwidth}rr}
\toprule
\textbf{Tablespace} & \textbf{Contenido} & \textbf{Estimación (300,000 empleados)} & \textbf{Asignado} & \textbf{Uso actual} \\
\midrule
\texttt{ts\_\allowbreak{}rrhh\_\allowbreak{}data} & Catálogos, historial, niveles MAC y tablas de afinación & $<$ 1\,MB & 32\,MB & 100\,kB \\
\texttt{ts\_\allowbreak{}rrhh\_\allowbreak{}part} & Particiones activas de \texttt{empleados} & 300,000 $\times$ 120\,B $\approx$ 35\,MB & 128\,MB & 4,452\,kB \\
\texttt{ts\_\allowbreak{}rrhh\_\allowbreak{}idx} & Llaves primarias e índices secundarios & 300,000 $\times$ 115\,B $\approx$ 33\,MB & 128\,MB & 6,772\,kB \\
\texttt{ts\_\allowbreak{}rrhh\_\allowbreak{}audit} & Bitácora de auditoría y registro del mantenimiento & 1,000,000 eventos $\times$ 120\,B $\approx$ 115\,MB & 256\,MB & 68\,kB \\
\texttt{ts\_\allowbreak{}rrhh\_\allowbreak{}hist} & Particiones históricas (antes de 2021) & 3\,MB hoy, crece al archivar & 64\,MB & 3,020\,kB \\
\bottomrule
\end{tabular}
\caption{Tamaño definido para cada tablespace. En los que crecen con los empleados (\texttt{part}, \texttt{idx}, \texttt{audit}) el asignado deja una holgura de 2 a 4 veces sobre la estimación; \texttt{data} e \texttt{hist} tienen un mínimo fijo. El uso actual es el medido al final del proyecto, después de eliminar el índice de control y de la primera corrida del mantenimiento programado.}
\label{tab:dimensionamiento}
\end{table}

\subsection{Roles, usuarios y acceso a la base}

Los privilegios se agrupan en tres roles sin inicio de sesión (\texttt{rol\_\allowbreak{}admin}, \texttt{rol\_\allowbreak{}analista}, \texttt{rol\_\allowbreak{}auditor}), y a cada integrante del equipo se le crea un usuario ficticio con contraseña y un rol: \texttt{usr\_\allowbreak{}admin}, \texttt{usr\_\allowbreak{}analista\_\allowbreak{}karla} y \texttt{usr\_\allowbreak{}auditor\_\allowbreak{}victor}. Un cuarto rol, \texttt{rol\_\allowbreak{}propietario}, es el dueño del esquema y de sus objetos; no inicia sesión ni tiene miembros, de modo que ningún usuario de trabajo es dueño de nada (sección~\ref{sec:seguridad}). Se retira el acceso por defecto a \texttt{PUBLIC} y solo el administrador puede crear objetos en los tablespaces.

\codigo{ev03_roles_usuarios}{Roles personalizados y usuarios del equipo}
\evid{ev03_roles_usuarios}{Creación de roles, usuarios y asignación de rol.}{ev03}

\codigo{ev04_esquema_acceso}{Control de acceso a la base, tablespaces, esquema \texttt{rrhh} y tamaño asignado a cada tablespace}
\evid{ev04_esquema_acceso}{Restricción de acceso a la base, listado de roles con \texttt{\textbackslash du} y tamaño asignado a cada tablespace (\texttt{capacidad\_\allowbreak{}tablespaces}).}{ev04}

\subsection{Esquema relacional y particionamiento}

El esquema \texttt{rrhh} (propiedad de \texttt{rol\_\allowbreak{}propietario}; \texttt{usr\_\allowbreak{}admin} lo construye y al cerrar el diseño la propiedad se transfiere) contiene:

\begin{itemize}
    \item \texttt{roles} y \texttt{departamentos}: catálogos en \texttt{ts\_\allowbreak{}rrhh\_\allowbreak{}data}, con sus índices de llave en \texttt{ts\_\allowbreak{}rrhh\_\allowbreak{}idx}.
    \item \texttt{empleados}: tabla principal \textbf{particionada por rango de \texttt{fecha\_\allowbreak{}ingreso}} en cinco particiones ubicadas en \texttt{ts\_\allowbreak{}rrhh\_\allowbreak{}part}. Incluye la columna \texttt{clasificacion\_\allowbreak{}nivel} (1 Público, 2 Confidencial, 3 Secreto) que alimenta el control MAC. La llave primaria es compuesta (\texttt{emp\_\allowbreak{}id}, \texttt{fecha\_\allowbreak{}ingreso}) porque PostgreSQL exige incluir la columna de partición.
    \item \texttt{historial\_\allowbreak{}salarios}: historial con llave foránea a la llave compuesta de \texttt{empleados}.
\end{itemize}

\codigo{ev05_catalogos}{Tablas de catálogo: roles y departamentos}
\evid{ev05_catalogos}{Creación de \texttt{roles} y \texttt{departamentos} desde la sesión de \texttt{usr\_\allowbreak{}admin}.}{ev05}

\codigo{ev06_empleados_particiones}{Tabla \texttt{empleados} particionada por rango}
\evid{ev06_empleados_particiones}{Tabla \texttt{empleados} y sus cinco particiones.}{ev06}

\codigo{ev07_historial_ubicacion}{Historial de salarios y ubicación física de cada objeto}
\evid{ev07_historial_ubicacion}{Historial de salarios y consulta de ubicación de tablas, particiones e índices por tablespace.}{ev07}

\subsection{Vista de negocio}

La vista \texttt{v\_\allowbreak{}resumen\_\allowbreak{}empleados} une empleados, departamentos y roles y expone solo las filas cuyo nivel de clasificación no excede la autorización del usuario que consulta (función \texttt{nivel\_\allowbreak{}usuario()}, definida con \texttt{SECURITY DEFINER} para leer la tabla \texttt{autorizaciones} sin exponerla). Es la vista del requerimiento y también el punto de acceso del analista.

La vista se crea con \texttt{security\_\allowbreak{}barrier}. Sin esa opción el optimizador puede evaluar antes un filtro escrito por el usuario que el filtro MAC, y los mensajes de error de ese filtro revelan datos de filas ocultas. Lo comprobamos con la versión anterior del proyecto: la analista obtenía \texttt{division by zero} con \texttt{1/(CASE WHEN clasificacion\_\allowbreak{}nivel = 3 THEN 0 ELSE 1 END)}, lo que prueba que existen filas secretas, y \texttt{invalid input syntax for type integer: "42085.07"} con \texttt{(salario::text)::int}, que muestra el salario de un empleado de nivel 3. Con \texttt{security\_\allowbreak{}barrier} el filtro MAC se aplica primero y la misma consulta devuelve las 40,000 filas visibles sin error (evidencia~\ref{fig:ev14}).

\codigo{ev08_mac_vista}{Niveles MAC, función \texttt{nivel\_\allowbreak{}usuario()} y vista \texttt{v\_\allowbreak{}resumen\_\allowbreak{}empleados}}
\evid{ev08_mac_vista}{Tabla de autorizaciones, función y vista de negocio.}{ev08}

% ==========================================
% 4. SEGURIDAD
% ==========================================
\section{Seguridad multinivel}
\label{sec:seguridad}

Los tres modelos se combinan en capas:

\begin{description}
    \item[RBAC (basado en roles):] los privilegios se otorgan a los roles, no a las personas; para dar o quitar acceso basta con cambiar la pertenencia a un rol.
    \item[DAC (discrecional):] quien tiene un privilegio con \texttt{WITH GRANT OPTION} decide con \texttt{GRANT}/\texttt{REVOKE} a quién se lo delega. El administrador recibe así la lectura de \texttt{empleados}, y se demuestra otorgando y retirando \texttt{SELECT} a Karla.
    \item[MAC (obligatorio):] aun con \texttt{GRANT}, el usuario solo ve las filas con \texttt{clasificacion\_\allowbreak{}nivel} menor o igual a su autorización, mediante políticas \textit{Row Level Security} con \texttt{FORCE}, que se aplican también a quien recibe permisos directos sobre la tabla. Niveles: administrador 3, Karla 2, Víctor 3. El MAC restringe al analista y al auditor; el administrador tiene acceso total por diseño (la política \texttt{pol\_\allowbreak{}admin\_\allowbreak{}escritura} le permite todas las filas), y el superusuario \texttt{postgres} siempre omite RLS, como ocurre con \texttt{SYS} en Oracle.
    \item[Separación de funciones:] en PostgreSQL el dueño de una tabla puede desactivar sus triggers, quitarle el RLS o borrarla, y el dueño de un esquema puede borrar cualquier objeto dentro de él. Por eso el administrador \textbf{no es dueño de nada}: \texttt{REASSIGN OWNED BY usr\_\allowbreak{}admin TO rol\_\allowbreak{}propietario} transfiere los objetos que creó, y \texttt{rol\_\allowbreak{}admin} solo recibe \texttt{SELECT}, \texttt{INSERT}, \texttt{UPDATE} y \texttt{DELETE}, sin \texttt{TRUNCATE} (vaciaría la tabla sin disparar la auditoría por fila) ni \texttt{TRIGGER}, y solo lectura sobre la bitácora y sobre \texttt{autorizaciones} (si pudiera escribir ahí, podría subirle el nivel MAC a cualquiera). En la versión anterior del proyecto comprobamos que el administrador podía desactivar el trigger de auditoría, borrar la bitácora con \texttt{DROP TABLE}, vaciar \texttt{empleados} con \texttt{TRUNCATE ... CASCADE} y elevar a Karla a nivel 3; ahora las cuatro operaciones se rechazan (evidencia~\ref{fig:ev15}).
\end{description}

\begin{table}[H]
\centering
\small
\renewcommand{\arraystretch}{1.3}
\begin{tabular}{lp{0.28\textwidth}p{0.38\textwidth}}
\toprule
\textbf{Rol} & \textbf{Usuario (nivel MAC)} & \textbf{Privilegios} \\
\midrule
\texttt{rol\_\allowbreak{}admin} & \texttt{usr\_\allowbreak{}admin} (3) & \texttt{SELECT}/\texttt{INSERT}/\texttt{UPDATE}/\texttt{DELETE} sobre los datos y \texttt{CREATE} en el esquema; solo \texttt{SELECT} en la bitácora y en \texttt{autorizaciones}; \texttt{SELECT} sobre \texttt{empleados} con \texttt{GRANT OPTION} \\
\texttt{rol\_\allowbreak{}propietario} & (ninguno) & Dueño del esquema y de sus objetos; sin inicio de sesión ni miembros \\
\texttt{rol\_\allowbreak{}analista} & \texttt{usr\_\allowbreak{}analista\_\allowbreak{}karla} (2) & \texttt{SELECT} sobre \texttt{v\_\allowbreak{}resumen\_\allowbreak{}empleados} y \texttt{departamentos} \\
\texttt{rol\_\allowbreak{}auditor} & \texttt{usr\_\allowbreak{}auditor\_\allowbreak{}victor} (3) & \texttt{SELECT} sobre \texttt{audit\_\allowbreak{}log\_\allowbreak{}empleados} y la vista \\
\bottomrule
\end{tabular}
\caption{Roles, usuarios y privilegios otorgados.}
\end{table}

\codigo{ev11_rbac_dac}{Privilegios por rol (RBAC/DAC)}
\evid{ev11_rbac_dac}{Transferencia de la propiedad a \texttt{rol\_\allowbreak{}propietario} y asignación de privilegios a los tres roles.}{ev11}

\codigo{ev12_mac_rls}{Niveles MAC, políticas \textit{Row Level Security} y resumen de privilegios}
\evid{ev12_mac_rls}{Políticas MAC, privilegios efectivos por rol (\texttt{information\_\allowbreak{}schema.role\_\allowbreak{}table\_\allowbreak{}grants}) y propietario de cada objeto.}{ev12}

% ==========================================
% 5. AUDITORIA
% ==========================================
\section{Auditoría y control automático de cambios}
\label{sec:auditoria}

La bitácora \texttt{audit\_\allowbreak{}log\_\allowbreak{}empleados} reside en su propio tablespace (\texttt{ts\_\allowbreak{}rrhh\_\allowbreak{}audit}). El trigger \texttt{trg\_\allowbreak{}audit\_\allowbreak{}empleados} se dispara después de cada \texttt{INSERT}, \texttt{UPDATE} o \texttt{DELETE} sobre \texttt{empleados} y registra la acción, el empleado, el usuario de sesión, la fecha, el salario anterior y el nuevo y la IP de origen. Además, cuando cambia el salario inserta una fila en \texttt{historial\_\allowbreak{}salarios}.

Dos decisiones de diseño:

\begin{itemize}
    \item La función del trigger es \texttt{SECURITY DEFINER}: registra el evento aunque el usuario que modifica no tenga ningún privilegio sobre la bitácora.
    \item La bitácora es \textbf{inmutable}: los triggers \texttt{trg\_\allowbreak{}bitacora\_\allowbreak{}inmutable} y \texttt{trg\_\allowbreak{}bitacora\_\allowbreak{}no\_\allowbreak{}truncate} bloquean \texttt{UPDATE}, \texttt{DELETE} y \texttt{TRUNCATE} para cualquier usuario que use SQL normal, incluido \texttt{postgres}. El administrador no tiene permiso de escritura sobre ella y, como no es dueño de los objetos, tampoco puede borrarla ni desactivar el trigger de auditoría de \texttt{empleados}. \textbf{Límite:} un superusuario podría desactivar antes los triggers (\texttt{ALTER TABLE ... DISABLE TRIGGER} o \texttt{session\_\allowbreak{}replication\_\allowbreak{}role = replica}, el mismo mecanismo que se usa en la carga masiva); por eso, en producción, la bitácora debe copiarse además a un servidor externo (sección~\ref{sec:recomendaciones}).
\end{itemize}

\codigo{ev09_trigger_auditoria}{Bitácora y trigger de auditoría de empleados}
\evid{ev09_trigger_auditoria}{Creación de la bitácora, la función de auditoría y el trigger.}{ev09}

\codigo{ev10_bitacora_inmutable}{Bitácora inmutable}
\evid{ev10_bitacora_inmutable}{Triggers que protegen la bitácora y listado de triggers del esquema (todos \texttt{ACTIVO}).}{ev10}

% ==========================================
% 6. CARGA MASIVA
% ==========================================
\section{Inserción masiva de datos con PL/pgSQL}
\label{sec:carga}

Se cargan \textbf{60,000 empleados simulados} (el mínimo exigido es 50,000) con un bloque anónimo \texttt{DO} escrito en PL/pgSQL, equivalente al bloque PL/SQL de Oracle. Explicación del código:

\begin{enumerate}
    \item Se insertan primero los catálogos (4 roles y 4 departamentos) para satisfacer las llaves foráneas.
    \item \texttt{session\_\allowbreak{}replication\_\allowbreak{}role = replica} desactiva los triggers durante la carga: sin esto se generarían 60,000 eventos artificiales en la bitácora. Las operaciones reales posteriores sí se auditan. Se restablece a \texttt{origin} al terminar.
    \item \texttt{setseed(0.42)} fija la semilla de \texttt{random()}, por lo que los salarios son reproducibles en cada ejecución.
    \item El bloque \texttt{DO} recorre \texttt{FOR v\_\allowbreak{}ini IN 1..v\_\allowbreak{}total BY v\_\allowbreak{}lote} e inserta lotes de 5,000 filas con un \texttt{INSERT ... SELECT} sobre \texttt{generate\_\allowbreak{}series}, mucho más eficiente que insertar fila por fila. Hace \texttt{COMMIT} por lote (permitido dentro de \texttt{DO}) y avisa el avance con \texttt{RAISE NOTICE}.
    \item Los datos se distribuyen de forma determinista: departamento $i \bmod 4$ y rol $\lfloor i/4 \rfloor \bmod 4$ (independientes entre sí, así que aparecen las 16 combinaciones departamento--rol con 3,750 empleados cada una), nivel de clasificación $1 + (i \bmod 3)$, salario aleatorio entre 12,000 y 95,000 y fechas de ingreso desde 2018-01-01 (cubren cuatro de las cinco particiones). \texttt{ANALYZE} actualiza las estadísticas.
\end{enumerate}

\codigo{ev13_carga_masiva}{Carga masiva de 60,000 empleados con un bloque \texttt{DO}}
\evid{ev13_carga_masiva}{Ejecución de la carga: 12 lotes, distribución por partición y total de 60,000 registros.}{ev13}

Distribución resultante por partición:

\begin{table}[H]
\centering
\small
\renewcommand{\arraystretch}{1.2}
\begin{tabular}{lrcc}
\toprule
\textbf{Partición} & \textbf{Registros} & \textbf{Desde} & \textbf{Hasta} \\
\midrule
\texttt{p\_\allowbreak{}historico} & 21,920 & 2018-01-01 & 2020-12-31 \\
\texttt{p\_\allowbreak{}2021\_\allowbreak{}2022} & 14,600 & 2021-01-01 & 2022-12-31 \\
\texttt{p\_\allowbreak{}2023\_\allowbreak{}2024} & 14,620 & 2023-01-01 & 2024-12-31 \\
\texttt{p\_\allowbreak{}2025\_\allowbreak{}2026} & 8,860 & 2025-01-01 & 2026-03-19 \\
\texttt{p\_\allowbreak{}futuro} & 0 & --- & --- \\
\midrule
\textbf{Total} & \textbf{60,000} & & \\
\bottomrule
\end{tabular}
\caption{Registros cargados por partición.}
\end{table}

% ==========================================
% 7. PRUEBAS DE SEGURIDAD
% ==========================================
\section{Pruebas de acceso, modificación y auditoría por rol}
\label{sec:pruebas}

Cada prueba se hace con una conexión TCP (\texttt{psql -h localhost}) con la contraseña del usuario, de modo que se valida también la autenticación. Los mensajes \texttt{ERROR} que aparecen en las capturas son \textbf{resultados esperados}: demuestran que el control de acceso funciona.

\subsection{Analista (\texttt{usr\_\allowbreak{}analista\_\allowbreak{}karla}, nivel 2)}

La primera conexión usa una contraseña incorrecta y es rechazada. Con la correcta, Karla ve 40,000 de las 60,000 filas a través de la vista (MAC) y puede consultar departamentos. El intento de descubrir filas ocultas con una división entre cero devuelve las mismas 40,000 filas sin error, gracias a \texttt{security\_\allowbreak{}barrier}. No puede leer \texttt{empleados}, borrar, leer la bitácora ni crear tablas.

\codigo{ev14_prueba_analista}{Pruebas del rol analista}
\evid{ev14_prueba_analista}{Sesión de Karla: contraseña incorrecta, lectura permitida a la vista y denegaciones.}{ev14}

\subsection{Administrador (\texttt{usr\_\allowbreak{}admin}, nivel 3)}

El administrador modifica un salario, inserta y elimina un empleado, y crea y borra una tabla. Se le niega modificar la bitácora, y también desactivar el trigger de auditoría, borrar la bitácora con \texttt{DROP TABLE}, vaciar \texttt{empleados} con \texttt{TRUNCATE} y cambiar los niveles MAC: no es dueño de los objetos (separación de funciones). El trigger registró las tres operaciones y el cambio de salario alimentó \texttt{historial\_\allowbreak{}salarios}.

\codigo{ev15_prueba_admin}{Pruebas del rol administrador}
\evid{ev15_prueba_admin}{Sesión del administrador: DML, DDL, bitácora protegida, intentos de evadir la auditoría o cambiar el MAC (rechazados) e historial de salarios.}{ev15}

\subsection{Auditor (\texttt{usr\_\allowbreak{}auditor\_\allowbreak{}victor}, nivel 3)}

Víctor lee la bitácora y la ve completa (tres eventos del administrador con IP \texttt{::1}), resume los eventos por usuario y acción, y ve las 60,000 filas de la vista; no puede modificar la bitácora ni borrar empleados.

\codigo{ev16_prueba_auditor}{Pruebas del rol auditor}
\evid{ev16_prueba_auditor}{Sesión del auditor: lectura de la bitácora y restricciones de escritura.}{ev16}

\subsection{DAC, MAC e inmutabilidad}

El administrador otorga \texttt{SELECT} sobre \texttt{empleados} a Karla (DAC): aun así solo ve 40,000 filas con nivel máximo 2, porque el MAC se mantiene. Tras el \texttt{REVOKE}, la consulta vuelve a ser denegada. Finalmente, ni siquiera \texttt{postgres} puede vaciar o borrar la bitácora con \texttt{TRUNCATE} o \texttt{DELETE}: los triggers lo impiden.

\codigo{ev17_dac_mac_inmutable}{DAC, MAC y bitácora inmutable}
\evid{ev17_dac_mac_inmutable}{\texttt{GRANT}/\texttt{REVOKE} discrecional, filtrado MAC y bloqueo de \texttt{TRUNCATE}/\texttt{DELETE} en la bitácora.}{ev17}

\begin{table}[H]
\centering
\small
\renewcommand{\arraystretch}{1.3}
\begin{tabular}{p{0.37\textwidth}ccc}
\toprule
\textbf{Operación} & \textbf{Admin} & \textbf{Analista} & \textbf{Auditor} \\
\midrule
\texttt{SELECT} en \texttt{v\_\allowbreak{}resumen\_\allowbreak{}empleados} & 60,000 filas & 40,000 filas & 60,000 filas \\
\texttt{SELECT} en \texttt{empleados} & Sí & Denegado & Denegado \\
\texttt{SELECT} en \texttt{departamentos} & Sí & Sí & Sin privilegio$^{*}$ \\
\texttt{INSERT}/\texttt{UPDATE}/\texttt{DELETE} en \texttt{empleados} & Sí & Denegado & Denegado \\
\texttt{SELECT} en \texttt{audit\_\allowbreak{}log\_\allowbreak{}empleados} & Sí & Denegado & Sí \\
\texttt{UPDATE} en \texttt{audit\_\allowbreak{}log\_\allowbreak{}empleados} & Denegado & Denegado & Denegado \\
\texttt{CREATE TABLE} en \texttt{rrhh} & Sí & Denegado & Sin privilegio$^{*}$ \\
\texttt{DISABLE TRIGGER} de auditoría / \texttt{DROP} de la bitácora & Denegado & Sin privilegio$^{*}$ & Sin privilegio$^{*}$ \\
\texttt{TRUNCATE} en \texttt{empleados} & Denegado & Sin privilegio$^{*}$ & Sin privilegio$^{*}$ \\
\texttt{UPDATE} en \texttt{autorizaciones} (nivel MAC) & Denegado & Sin privilegio$^{*}$ & Sin privilegio$^{*}$ \\
\bottomrule
\end{tabular}
\caption{Matriz de resultados de las pruebas de acceso. $^{*}$El rol no recibió ese privilegio (sección~\ref{sec:seguridad}); la operación no se ejecutó en sus pruebas.}
\label{tab:matriz}
\end{table}

% ==========================================
% 8. RENDIMIENTO: ANTES
% ==========================================
\section{Optimización y afinación del rendimiento}

\subsection{Infraestructura de medición}

Para que la comparación sea objetiva y repetible se guardan las tres consultas complejas en la tabla \texttt{consultas\_\allowbreak{}tuning} y se crean funciones PL/pgSQL que ejecutan \texttt{EXPLAIN (FORMAT JSON)} (costo estimado del optimizador, análogo al costo de \texttt{DBMS\_\allowbreak{}XPLAN}) y \texttt{EXPLAIN (ANALYZE)} (tiempo real de ejecución). La función \texttt{registrar\_\allowbreak{}etapa} guarda en \texttt{tuning\_\allowbreak{}resultados} el costo y el \textbf{mejor tiempo de cinco ejecuciones} de cada consulta en cada etapa, para reducir el ruido de una sola medición. La función \texttt{ejecutar\_\allowbreak{}carga(n)} repite las consultas $n$ veces y sirve para generar la carga de trabajo que analiza el advisor.

\begin{itemize}
    \item \textbf{Q1:} salario promedio por departamento en un rango de fechas (poda de particiones).
    \item \textbf{Q2:} personal y gasto de nómina por departamento y rol con \texttt{JOIN} a ambos catálogos, para salarios superiores a 92,000.
    \item \textbf{Q3:} empleados del departamento 10 con salario mayor a 1.7 veces el promedio del departamento (subconsulta).
\end{itemize}

\codigo{ev18_consultas_tuning}{Consultas a analizar y tabla de resultados por etapa}
\evid{ev18_consultas_tuning}{Registro de las tres consultas complejas.}{ev18}

\codigo{ev19_funciones_plan}{Funciones de medición, tabla de resultados y carga de trabajo}
\evid{ev19_funciones_plan}{Funciones \texttt{costo\_\allowbreak{}plan}, \texttt{tiempo\_\allowbreak{}plan}, \texttt{registrar\_\allowbreak{}etapa} y \texttt{ejecutar\_\allowbreak{}carga}.}{ev19}

\subsection{Medición inicial (sin índices secundarios)}
\label{sec:antes}

Tras \texttt{VACUUM ANALYZE} se mide la etapa base. Solo existen las llaves primarias.

\codigo{ev20_medicion_base}{Medición de la etapa base}
\evid{ev20_medicion_base}{Costo y tiempo de las tres consultas sin índices secundarios.}{ev20}

\paragraph{Q1.} El optimizador poda las particiones y solo lee \texttt{p\_\allowbreak{}2023\_\allowbreak{}2024}, pero lo hace con un \textit{Seq Scan} completo de la partición (costo 432.81).

\codigo{ev21_plan_q1_antes}{\texttt{EXPLAIN} de Q1 antes de los índices}
\evid{ev21_plan_q1_antes}{Plan de Q1 antes de indexar: \textit{Seq Scan} en una partición.}{ev21}

\paragraph{Q2.} Recorre secuencialmente las cuatro particiones con datos, filtra por salario y hace dos \textit{Hash Join} con los catálogos (costo 1,708.78).

\codigo{ev22_plan_q2_antes}{\texttt{EXPLAIN} de Q2 antes de los índices}
\evid{ev22_plan_q2_antes}{Plan de Q2 antes de indexar.}{ev22}

\paragraph{Q3.} La subconsulta y la consulta externa recorren dos veces todas las particiones (costo 3,497.53, el mayor de las tres).

\codigo{ev23_plan_q3_antes}{\texttt{EXPLAIN} de Q3 antes de los índices}
\evid{ev23_plan_q3_antes}{Plan de Q3 antes de indexar.}{ev23}

% ==========================================
% 9. ADVISOR
% ==========================================
\subsection{Advisor: hallazgos y recomendaciones}
\label{sec:advisor}

PostgreSQL no incluye un asesor automático equivalente a \texttt{DBA\_\allowbreak{}ADVISOR\_\allowbreak{}FINDINGS} y \texttt{DBA\_\allowbreak{}ADVISOR\_\allowbreak{}RECOMMENDATIONS}. Se construyó uno con tres vistas sobre las estadísticas del motor:

\begin{itemize}
    \item \texttt{v\_\allowbreak{}advisor\_\allowbreak{}hallazgos} (análoga a \texttt{DBA\_\allowbreak{}ADVISOR\_\allowbreak{}FINDINGS}): particiones con lecturas secuenciales frecuentes (\texttt{seq\_\allowbreak{}scan} mayor al doble de \texttt{idx\_\allowbreak{}scan}) e índices que nunca se usan.
    \item \texttt{v\_\allowbreak{}advisor\_\allowbreak{}recomendaciones} (análoga a \texttt{DBA\_\allowbreak{}ADVISOR\_\allowbreak{}RECOMMENDATIONS}): acción sugerida por tipo de hallazgo, con el número de objetos afectados.
    \item \texttt{v\_\allowbreak{}advisor\_\allowbreak{}sentencias}: sentencias más costosas de \texttt{pg\_\allowbreak{}stat\_\allowbreak{}statements}, de donde se toman las columnas por indexar.
\end{itemize}

\codigo{ev24_advisor_vistas}{Vistas del advisor y ejecución de la carga de trabajo}
\evid{ev24_advisor_vistas}{Creación de las vistas del advisor.}{ev24}

Con la carga de trabajo ejecutada 20 veces por consulta, el advisor reporta \textbf{4 particiones con lecturas secuenciales} (5,285,500 filas leídas en total) y recomienda crear índices sobre las columnas de filtro de las sentencias más costosas.

\evid{ev25_advisor_hallazgos}{Hallazgos y recomendación del advisor antes de indexar.}{ev25}

\begin{table}[H]
\centering
\small
\renewcommand{\arraystretch}{1.2}
\begin{tabular}{lrr}
\toprule
\textbf{Partición} & \textbf{Lecturas secuenciales} & \textbf{Filas leídas} \\
\midrule
\texttt{p\_\allowbreak{}historico} & 83 & 1,797,440 \\
\texttt{p\_\allowbreak{}2023\_\allowbreak{}2024} & 108 & 1,564,340 \\
\texttt{p\_\allowbreak{}2021\_\allowbreak{}2022} & 83 & 1,197,200 \\
\texttt{p\_\allowbreak{}2025\_\allowbreak{}2026} & 83 & 726,520 \\
\bottomrule
\end{tabular}
\caption{Hallazgo del advisor: \texttt{LECTURAS SECUENCIALES}.}
\end{table}

% ==========================================
% 10. INDICES
% ==========================================
\subsection{Creación y recreación de índices}
\label{sec:indices}

Siguiendo las recomendaciones se crean cuatro índices sobre \texttt{empleados} (PostgreSQL crea uno local por partición, en \texttt{ts\_\allowbreak{}rrhh\_\allowbreak{}idx}): \texttt{idx\_\allowbreak{}emp\_\allowbreak{}fecha\_\allowbreak{}depto} para Q1, \texttt{idx\_\allowbreak{}emp\_\allowbreak{}depto\_\allowbreak{}salario} para Q3, \texttt{idx\_\allowbreak{}emp\_\allowbreak{}salario} (con \texttt{INCLUDE}) para Q2 e \texttt{idx\_\allowbreak{}emp\_\allowbreak{}email}. Este último es un \textbf{índice de control}: ninguna consulta de la carga filtra por \texttt{email}, y se crea a propósito para comprobar que el advisor detecta índices sin uso (sección~\ref{sec:aplicacion}). Se mide la etapa \texttt{2\_\allowbreak{}con\_\allowbreak{}indices} y se consulta el tamaño de cada índice.

\codigo{ev26_crear_indices}{Creación de índices y medición}
\evid{ev26_crear_indices}{Creación de índices, medición de la etapa 2 y tamaño de cada índice.}{ev26}

\paragraph{Recreación de un índice.} Con \texttt{idx\_\allowbreak{}emp\_\allowbreak{}fecha\_\allowbreak{}depto}, Q1 baja de 432.81 a 232.80, pero aún debe ir a la tabla a leer \texttt{salario}. Se elimina y se recrea como \textbf{índice de cobertura} (\texttt{INCLUDE (depto\_\allowbreak{}id, salario)}), con lo que el plan pasa a \textit{Index Only Scan}: el costo cae a 28.68.

\codigo{ev27_recrear_indice}{Recreación del índice de Q1 como índice de cobertura}
\evid{ev27_recrear_indice}{Eliminación y recreación del índice y medición de la etapa 3.}{ev27}

\subsubsection{Planes después de la afinación}

\codigo{ev28_plan_q1_despues}{\texttt{EXPLAIN} de Q1 después de los índices}
\evid{ev28_plan_q1_despues}{Q1 usa \textit{Index Only Scan} sobre el índice de cobertura (costo 28.68).}{ev28}

\codigo{ev29_plan_q2_despues}{\texttt{EXPLAIN} de Q2 después de los índices}
\evid{ev29_plan_q2_despues}{Q2 sustituye los \textit{Seq Scan} por \textit{Bitmap Index Scan} (costo 1,053.14).}{ev29}

\codigo{ev30_plan_q3_despues}{\texttt{EXPLAIN} de Q3 después de los índices}
\evid{ev30_plan_q3_despues}{Q3 usa \textit{Index Only Scan} en la subconsulta y \textit{Bitmap Scan} en la consulta externa (costo 1,719.56).}{ev30}

\codigo{ev31_comparativo}{Comparativo del costo y tiempo por etapa}
\evid{ev31_comparativo}{Comparativo de costo y tiempo antes y después de indexar.}{ev31}

% ==========================================
% 11. PARAMETROS
% ==========================================
\subsection{Ajuste de parámetros de configuración}
\label{sec:parametros}

Se consultan los parámetros críticos: \texttt{block\_\allowbreak{}size} (8192 bytes, equivalente a \texttt{DB\_\allowbreak{}BLOCK\_\allowbreak{}SIZE}), \texttt{shared\_\allowbreak{}buffers} (16384 bloques de 8\,kB = 128\,MB, análogo al SGA), \texttt{work\_\allowbreak{}mem}, \texttt{maintenance\_\allowbreak{}work\_\allowbreak{}mem}, \texttt{effective\_\allowbreak{}cache\_\allowbreak{}size}, \texttt{random\_\allowbreak{}page\_\allowbreak{}cost} y \texttt{max\_\allowbreak{}connections}.

\codigo{ev37_parametros_actuales}{Valores actuales de los parámetros críticos}
\evid{ev37_parametros_actuales}{Valores iniciales de los parámetros; la columna \texttt{context} indica cuáles requieren reinicio.}{ev37}

\textbf{\texttt{block\_\allowbreak{}size} no es modificable}: su contexto es \texttt{internal}, es decir, se fija al compilar el motor (en Oracle \texttt{DB\_\allowbreak{}BLOCK\_\allowbreak{}SIZE} se fija al crear la base), por lo que no se puede cambiar ni en caliente ni con reinicio. Los parámetros ajustados con \texttt{ALTER SYSTEM}, que se aplican con \texttt{pg\_\allowbreak{}reload\_\allowbreak{}conf()} sin reiniciar, fueron:

\begin{table}[H]
\centering
\small
\renewcommand{\arraystretch}{1.3}
\begin{tabular}{llp{0.42\textwidth}}
\toprule
\textbf{Parámetro} & \textbf{Antes $\rightarrow$ después} & \textbf{Justificación} \\
\midrule
\texttt{work\_\allowbreak{}mem} & 4\,MB $\rightarrow$ 32\,MB & Ordenamientos y \textit{Hash Join} sin desbordar a disco \\
\texttt{maintenance\_\allowbreak{}work\_\allowbreak{}mem} & 64\,MB $\rightarrow$ 256\,MB & \texttt{REINDEX}, \texttt{CREATE INDEX} y \texttt{VACUUM} más rápidos \\
\texttt{random\_\allowbreak{}page\_\allowbreak{}cost} & 4 $\rightarrow$ 1.1 & Almacenamiento SSD: el acceso por índice es más barato que lo que supone el valor por defecto \\
\bottomrule
\end{tabular}
\caption{Parámetros ajustados.}
\end{table}

\codigo{ev38_alter_system}{Ajuste de parámetros con \texttt{ALTER SYSTEM}}
\evid{ev38_alter_system}{Ajuste y verificación de los nuevos valores (fuente: \texttt{configuration file}).}{ev38}

Al repetir la medición (etapa 4) los costos estimados bajan de nuevo, sobre todo en Q1 y Q3:

\codigo{ev39_reevaluacion}{Reevaluación con los parámetros ajustados}
\evid{ev39_reevaluacion}{Costo y tiempo antes y después del ajuste de parámetros.}{ev39}

\subsection{Aplicación de la recomendación del advisor}
\label{sec:aplicacion}

Ninguna de las tres consultas filtra por \texttt{email}, de modo que \texttt{idx\_\allowbreak{}emp\_\allowbreak{}email} solo agrega costo de escritura y 3,000\,kB de espacio (el mayor de los índices). Se aplica la recomendación de eliminarlo. Después se reinician las estadísticas, se repite la carga de trabajo y se vuelve a consultar el advisor.

\codigo{ev40_aplicar_recomendacion}{Eliminación del índice sin beneficio de lectura}
\evid{ev40_aplicar_recomendacion}{Aplicación de la recomendación: \texttt{DROP INDEX}.}{ev40}

\codigo{ev41_verificacion_advisor}{Verificación del advisor después de la afinación}
\evid{ev41_verificacion_advisor}{Las cuatro particiones con datos ya no tienen lecturas secuenciales (\texttt{seq\_\allowbreak{}scan} = 0, \texttt{idx\_\allowbreak{}scan} entre 60 y 80).}{ev41}

El hallazgo \texttt{LECTURAS SECUENCIALES} desapareció: las particiones \texttt{p\_\allowbreak{}historico}, \texttt{p\_\allowbreak{}2021\_\allowbreak{}2022}, \texttt{p\_\allowbreak{}2023\_\allowbreak{}2024} y \texttt{p\_\allowbreak{}2025\_\allowbreak{}2026} ya solo se leen por índice. Las únicas lecturas secuenciales restantes son de las particiones vacías (\texttt{p\_\allowbreak{}2027\_\allowbreak{}2028}, \texttt{p\_\allowbreak{}futuro}), donde el optimizador prefiere \textit{Seq Scan} porque no tienen filas. El advisor reporta ahora ocho hallazgos \texttt{INDICE SIN USO}, que corresponden a índices de particiones que la carga de trabajo no toca (Q1 solo consulta \texttt{p\_\allowbreak{}2023\_\allowbreak{}2024}) o que están vacías; se documentan y no se eliminan porque son parte de la definición lógica de un índice particionado, que se mantiene para futuras consultas.

% ==========================================
% 12. COMPARATIVO
% ==========================================
\section{Análisis comparativo antes/después}
\label{sec:comparativo}

\begin{table}[H]
\centering
\small
\renewcommand{\arraystretch}{1.3}
\setlength{\tabcolsep}{4pt}
\begin{tabular}{lrrrrr}
\toprule
\textbf{Consulta} & \textbf{Sin índ.} & \textbf{Con índ.} & \textbf{Recreado} & \textbf{+ Parám.} & \textbf{Reduc.} \\
\midrule
Q1 (fechas) & 432.81 & 232.80 & 28.68 & 19.98 & 95.4\,\% \\
Q2 (\texttt{JOIN}) & 1,708.78 & 1,053.14 & 1,053.14 & 1,031.86 & 39.6\,\% \\
Q3 (subcons.) & 3,497.53 & 1,719.56 & 1,719.56 & 1,480.70 & 57.7\,\% \\
\bottomrule
\end{tabular}
\caption{Costo estimado del optimizador (unidades de costo de PostgreSQL) en cada etapa de la afinación.}
\end{table}

\begin{table}[H]
\centering
\small
\renewcommand{\arraystretch}{1.3}
\begin{tabular}{lrrrr}
\toprule
\textbf{Consulta} & \textbf{Sin índ. (ms)} & \textbf{Recreado (ms)} & \textbf{+ Parám. (ms)} & \textbf{Reduc.} \\
\midrule
Q1 & 0.566 & 0.096 & 0.099 & 82.5\,\% \\
Q2 & 5.051 & 1.287 & 1.341 & 73.5\,\% \\
Q3 & 9.231 & 2.538 & 2.497 & 72.9\,\% \\
\bottomrule
\end{tabular}
\caption{Tiempo real de ejecución (mejor de cinco ejecuciones con \texttt{EXPLAIN ANALYZE}).}
\end{table}

\begin{itemize}
    \item \textbf{Q1} es la que más mejora porque el filtro es muy selectivo: el índice de cobertura permite un \textit{Index Only Scan} sin tocar la tabla. Los índices sin cobertura (232.80) ya eran una mejora, pero recrearlo con \texttt{INCLUDE} multiplica el beneficio por ocho.
    \item \textbf{Q2} mejora menos porque debe leer cerca del 4\,\% de la tabla (2,234 de 60,000 filas estimadas) en todas las particiones y agregar; el índice evita leer todas las páginas, pero el costo del \textit{Hash Join} y la agregación se mantiene.
    \item \textbf{Q3} se beneficia en las dos partes de la consulta (subconsulta y filtro externo) y es la que más gana con el ajuste de \texttt{random\_\allowbreak{}page\_\allowbreak{}cost}.
    \item Los \textbf{costos} son estimaciones y son estables entre ejecuciones; los \textbf{milisegundos} varían ligeramente de una ejecución a otra. El ajuste de parámetros mejora el costo estimado pero, con este volumen (60,000 filas, que caben en memoria), prácticamente no cambia el tiempo real.
    \item \textbf{Costo de los índices:} los índices secundarios creados ocuparon 1,920\,kB (\texttt{idx\_\allowbreak{}emp\_\allowbreak{}depto\_\allowbreak{}salario}), 1,920\,kB (\texttt{idx\_\allowbreak{}emp\_\allowbreak{}salario}), 536\,kB (fecha) y 3,000\,kB (\texttt{idx\_\allowbreak{}emp\_\allowbreak{}email}), además del sobrecosto de escritura en cada \texttt{INSERT}/\texttt{UPDATE}; por eso se eliminó \texttt{idx\_\allowbreak{}emp\_\allowbreak{}email}, que no aportaba lecturas.
\end{itemize}

% ==========================================
% 13. ALMACENAMIENTO
% ==========================================
\section{Gestión avanzada de almacenamiento}
\label{sec:almacenamiento}

PostgreSQL no tiene \textit{datafiles} ni \texttt{AUTOEXTEND}: los archivos de un tablespace crecen mientras haya espacio en el volumen donde está su directorio. Por eso el control de crecimiento se hace contra el tamaño asignado a cada tablespace (tabla~\ref{tab:dimensionamiento}), con las vistas de uso, con la creación de nuevas particiones y con la reubicación de datos entre tablespaces.

\subsection{Monitoreo de tablespaces y particiones}

La vista \texttt{v\_\allowbreak{}uso\_\allowbreak{}tablespaces} (análoga a \texttt{DBA\_\allowbreak{}DATA\_\allowbreak{}FILES}/\texttt{DBA\_\allowbreak{}FREE\_\allowbreak{}SPACE}) muestra la ruta, el espacio usado, el tamaño asignado (tabla \texttt{capacidad\_\allowbreak{}tablespaces}), el porcentaje de uso, el estado (\texttt{OK} o \texttt{ALERTA} al llegar al 80\,\%) y el número de objetos de cada tablespace. \texttt{v\_\allowbreak{}particiones} muestra el rango, las filas y el tamaño de cada partición.

\codigo{ev32_uso_tablespaces}{Vista de uso de tablespaces frente al tamaño asignado}
\evid{ev32_uso_tablespaces}{Tamaño asignado y uso por tablespace: \texttt{ts\_\allowbreak{}rrhh\_\allowbreak{}idx} 10,236\,kB (7.8\,\% de 128\,MB) y \texttt{ts\_\allowbreak{}rrhh\_\allowbreak{}part} 6,972\,kB (5.3\,\%); todos en \texttt{OK}.}{ev32}

\codigo{ev33_particiones}{Vista de particiones}
\evid{ev33_particiones}{Rango, filas y tamaño de cada partición.}{ev33}

\subsection{Ajuste dinámico de particiones y tablespaces}

Sin detener el servicio ni perder datos se hacen dos ajustes: (1) se crea la nueva partición \texttt{p\_\allowbreak{}2027\_\allowbreak{}2028} desprendiendo \texttt{p\_\allowbreak{}futuro}, adjuntando la nueva y volviendo a adjuntar \texttt{p\_\allowbreak{}futuro} desde 2029; (2) se mueve la partición histórica y su índice al tablespace \texttt{ts\_\allowbreak{}rrhh\_\allowbreak{}hist} con \texttt{ALTER TABLE ... SET TABLESPACE}. Resultado: \texttt{ts\_\allowbreak{}rrhh\_\allowbreak{}hist} pasa de 0 a 3,020\,kB, \texttt{ts\_\allowbreak{}rrhh\_\allowbreak{}part} de 6,972 a 4,452\,kB y \texttt{ts\_\allowbreak{}rrhh\_\allowbreak{}idx} de 10,236 a 9,780\,kB, y el total sigue siendo de 60,000 registros.

\codigo{ev34_ajuste_dinamico}{Nueva partición y reubicación de datos históricos}
\evid{ev34_ajuste_dinamico}{Particiones y uso de tablespaces después del ajuste.}{ev34}

\subsection{Mantenimiento automático}

El procedimiento \texttt{prc\_\allowbreak{}mantenimiento(p\_\allowbreak{}prefijo)} recorre con un cursor implícito los índices secundarios cuyo nombre empieza con el prefijo y los reconstruye con \texttt{REINDEX INDEX} (SQL dinámico con \texttt{format('\%I')} para citar los identificadores de forma segura), y luego actualiza las estadísticas de todas las particiones con \texttt{ANALYZE}. Es el análogo de los scripts de \texttt{ALTER INDEX ... REBUILD} y \texttt{DBMS\_\allowbreak{}STATS}. Su ejecución periódica la hace el servicio \texttt{mantenimiento} descrito a continuación.

\codigo{ev35_mantenimiento}{Procedimiento de mantenimiento automático}
\evid{ev35_mantenimiento}{Creación y ejecución de \texttt{prc\_\allowbreak{}mantenimiento('p\_\allowbreak{}2025')}: índices reconstruidos y particiones analizadas.}{ev35}

\subsubsection{Programación automática del mantenimiento}
\label{sec:mant_programado}

Para que el mantenimiento no dependa de ejecutarlo a mano, \texttt{docker-compose\allowbreak{}.yaml} incluye un tercer servicio, \texttt{mantenimiento} (contenedor \texttt{bd-lab-final-mant}), análogo a un \textit{job} de \texttt{DBMS\_\allowbreak{}SCHEDULER}. Arranca cuando la base está lista, ejecuta \texttt{prc\_\allowbreak{}mantenimiento('p\_')} (todos los índices de las particiones y \texttt{ANALYZE}) y después repite cada \texttt{MANT\_\allowbreak{}INTERVALO\_\allowbreak{}SEG} segundos (86,400 = una vez al día). Cada corrida queda registrada en la tabla \texttt{log\_\allowbreak{}mantenimiento}, ubicada en el tablespace de auditoría y legible por el administrador y el auditor. Si los scripts del proyecto aún no se han ejecutado, la corrida se omite sin error.

\lstinputlisting[language={},caption={Servicio \texttt{mantenimiento} en docker-compose.yaml}, firstline=54, lastline=75]{../Dockerpfinal/docker-compose.yaml}

\lstinputlisting[language={},caption={Programador \texttt{mantenimiento/mantenimiento.sh}}]{../Dockerpfinal/mantenimiento/mantenimiento.sh}

\lstinputlisting[caption={Corrida registrada en \texttt{rrhh.log\_\allowbreak{}mantenimiento}}]{../Dockerpfinal/mantenimiento/mantenimiento_programado.sql}

Para la demostración se puede forzar una corrida inmediata y revisar su registro:

\begin{lstlisting}[language={},caption={Ejecución manual y verificación del mantenimiento programado}]
docker exec bd-lab-final-mant sh /mant/mantenimiento.sh --una-vez
docker logs bd-lab-final-mant
docker exec -it bd-lab-final psql -c "SELECT * FROM rrhh.log_mantenimiento;"
\end{lstlisting}

La evidencia muestra el registro del servicio, con la corrida que se omitió al arrancar porque los scripts aún no existían, y una corrida manual con su registro en \texttt{log\_\allowbreak{}mantenimiento}. También se puede demostrar en vivo con los comandos anteriores.

\evid{ev42_mant_programado}{Servicio \texttt{mantenimiento}: corrida omitida al arrancar y corrida manual registrada en \texttt{rrhh.log\_\allowbreak{}mantenimiento}.}{ev42}

\subsection{Truncado controlado de una partición}

Se cargan 1,000 filas de prueba en \texttt{p\_\allowbreak{}2027\_\allowbreak{}2028} y se vacía solo esa partición. Como una partición referenciada por una llave foránea (\texttt{historial\_\allowbreak{}salarios}) no admite \texttt{TRUNCATE} directo, se aplica la secuencia \texttt{DETACH PARTITION}, \texttt{TRUNCATE}, \texttt{ATTACH PARTITION} con el mismo rango. La partición pasa de 1,000 a 0 filas y los 60,000 empleados originales quedan intactos.

\codigo{ev36_truncado}{Truncado de una partición}
\evid{ev36_truncado}{Filas antes (1,000) y después (0) del truncado, y total de empleados (60,000).}{ev36}

% ==========================================
% 14. PROPUESTA
% ==========================================
\clearpage
\section{Propuesta para ambientes productivos}
\label{sec:recomendaciones}

Para llevar este diseño a producción proponemos lo siguiente:

\begin{itemize}
    \item \textbf{Seguridad:} aplicar siempre mínimo privilegio con \textit{RBAC}; no usar el superusuario en las aplicaciones; exigir \texttt{scram-sha-256} en todo acceso TCP, no dejar contraseñas en texto claro en los scripts (usar un gestor de secretos) y cifrar las conexiones con TLS.
    \item \textbf{Auditoría:} mantener la bitácora inmutable y en un tablespace propio, instalar \texttt{pgAudit} y enviar la bitácora a un servidor de registros externo. Así se cierra el único hueco que queda (un superusuario puede desactivar los triggers) y ni un administrador del sistema operativo puede alterarla.
    \item \textbf{Índices:} crear índices a partir de la evidencia de las vistas de estadísticas (no por intuición), preferir índices de cobertura para consultas frecuentes y revisar periódicamente los índices sin uso; cada índice cuesta espacio y escritura. Revisar cada semana \texttt{v\_\allowbreak{}advisor\_\allowbreak{}hallazgos} con la carga real y medir los planes también como analista y auditor, cuyas consultas pasan por la vista y por el RLS.
    \item \textbf{Particionamiento:} programar con \texttt{pg\_\allowbreak{}cron} la creación de la partición del año siguiente antes de que los datos lleguen a \texttt{p\_\allowbreak{}futuro}, y el traslado a \texttt{ts\_\allowbreak{}rrhh\_\allowbreak{}hist} de las particiones que cumplan cinco años, reutilizando \texttt{prc\_\allowbreak{}mantenimiento}.
    \item \textbf{Almacenamiento:} montar cada tablespace en un volumen del tamaño asignado en la tabla~\ref{tab:dimensionamiento} (no existe \texttt{AUTOEXTEND}), conectar \texttt{v\_\allowbreak{}uso\_\allowbreak{}tablespaces} a una alerta cuando el estado pase a \texttt{ALERTA} (80\,\%) y separar índices, datos, bitácora e históricos en discos distintos.
    \item \textbf{Parámetros:} dimensionar \texttt{shared\_\allowbreak{}buffers} (aprox. 25\,\% de la RAM), \texttt{work\_\allowbreak{}mem} según el número de conexiones concurrentes, y \texttt{random\_\allowbreak{}page\_\allowbreak{}cost} según el tipo de disco; tener presente que \texttt{block\_\allowbreak{}size} se decide al compilar el motor.
    \item \textbf{Operación:} tomar respaldos con \texttt{pg\_\allowbreak{}dump} o \texttt{pg\_\allowbreak{}basebackup} antes de cualquier cambio de estructura y ejecutar \texttt{VACUUM ANALYZE} después de cargas masivas, como se hizo en este proyecto.
\end{itemize}

% ==========================================
% 15. GUIA DEMO
% ==========================================
\section{Guía de demostración en clase (cambio de rol)}
\label{sec:demo}

Con los contenedores levantados (\texttt{docker compose start} desde \texttt{Dockerpfinal} si estaban detenidos) se pueden repetir en vivo las acciones más relevantes. Cada usuario se conecta con \texttt{docker exec -it bd-lab-final psql -h localhost -U <usuario> -d rrhh}:

\begin{enumerate}
    \item \textbf{Analista} (\texttt{usr\_\allowbreak{}analista\_\allowbreak{}karla}): \texttt{SELECT count(*) FROM v\_\allowbreak{}resumen\_\allowbreak{}empleados;} (40,000 filas); \texttt{SELECT * FROM empleados;} (denegado); \texttt{SELECT * FROM audit\_\allowbreak{}log\_\allowbreak{}empleados;} (denegado); el filtro con división entre cero \texttt{WHERE 1/(CASE WHEN clasificacion\_\allowbreak{}nivel = 3 THEN 0 ELSE 1 END) = 1} devuelve las mismas 40,000 filas sin error (\texttt{security\_\allowbreak{}barrier}).
    \item \textbf{Administrador} (\texttt{usr\_\allowbreak{}admin}): \texttt{UPDATE empleados SET salario = salario * 1.10 WHERE emp\_\allowbreak{}id = 100;}; \texttt{UPDATE audit\_\allowbreak{}log\_\allowbreak{}empleados ...} (denegado); \texttt{ALTER TABLE empleados DISABLE TRIGGER trg\_\allowbreak{}audit\_\allowbreak{}empleados;} (denegado: no es dueño de la tabla).
    \item \textbf{Auditor} (\texttt{usr\_\allowbreak{}auditor\_\allowbreak{}victor}): \texttt{SELECT * FROM audit\_\allowbreak{}log\_\allowbreak{}empleados;} muestra quién, qué y cuándo cambió el salario del empleado 100.
    \item \textbf{Superusuario} (\texttt{psql -U postgres -d rrhh}): \texttt{SELECT * FROM rrhh.v\_\allowbreak{}advisor\_\allowbreak{}hallazgos;}, \texttt{SELECT * FROM rrhh.v\_\allowbreak{}advisor\_\allowbreak{}recomendaciones;}, \texttt{SELECT * FROM rrhh.v\_\allowbreak{}uso\_\allowbreak{}tablespaces;}, \texttt{EXPLAIN} de las consultas Q1 a Q3 y \texttt{CALL rrhh.prc\_\allowbreak{}mantenimiento('p\_\allowbreak{}2025');}.
    \item \textbf{Cambio de rol sin reconectar} (desde \texttt{postgres}): \texttt{SET ROLE rol\_\allowbreak{}analista;} y \texttt{SET ROLE rol\_\allowbreak{}auditor;} seguidos de \texttt{RESET ROLE;} permiten mostrar los privilegios de cada rol. Para el nivel MAC, que depende del usuario de sesión (\texttt{session\_\allowbreak{}user}), se reconecta con cada usuario.
    \item \textbf{Mantenimiento automático:} \texttt{docker exec bd-lab-final-mant} \texttt{sh /mant/\allowbreak{}mantenimiento.sh} \texttt{-{}-una-vez} fuerza una corrida, y \texttt{SELECT * FROM rrhh.log\_\allowbreak{}mantenimiento;} muestra el registro de todas las corridas, incluidas las programadas.
\end{enumerate}

Para restaurar el escenario completo desde cero: \texttt{docker compose down -v} y ejecutar \texttt{python generar\_\allowbreak{}evidencias.py} (o los diez scripts de \texttt{sql/} en orden). El script \texttt{validar.sh} comprueba el estado final con 14 verificaciones.

% ==========================================
% 16. CIERRE
% ==========================================
\section{Cierre del entorno}

Al terminar se detienen los contenedores con \texttt{docker compose stop}, sin borrar los volúmenes, con lo que se conservan los datos del proyecto:

\evid{ev43_cierre}{Detención de los contenedores del proyecto y su estado final.}{ev43}

% ==========================================
% CONCLUSIONES
% ==========================================
\section{Conclusiones}
\label{sec:conclusiones}

\subsection{Actividades logradas}

\begin{itemize}
    \item \textbf{Diseño:} esquema completo de empleados, departamentos y roles, con \texttt{empleados} particionada por rango en cinco particiones y cinco tablespaces especializados (datos, índices, particiones, auditoría e histórico), cada uno con un tamaño definido por cálculo para 300,000 empleados y vigilado contra un umbral de 80\,\% (tabla~\ref{tab:dimensionamiento}).
    \item \textbf{Seguridad:} tres usuarios y tres roles de trabajo, más un rol propietario sin inicio de sesión (separación de funciones: el administrador no puede borrar la bitácora, desactivar la auditoría ni cambiar los niveles MAC); RBAC con privilegios por rol, DAC demostrado con \texttt{GRANT}/\texttt{REVOKE} y MAC con niveles de autorización, \textit{Row Level Security} y una vista \texttt{security\_\allowbreak{}barrier} que no deja ver filas ocultas a través de mensajes de error (el analista ve 40,000 de 60,000 filas). Las pruebas por rol (tabla~\ref{tab:matriz}) validan acceso, modificación y auditoría, incluso con rechazo de contraseña incorrecta.
    \item \textbf{Auditoría:} vista y trigger que registran usuario, IP y salarios anterior y nuevo, y alimentan el historial de salarios; la bitácora no puede modificarse con SQL normal, ni siquiera por \texttt{postgres}, mientras sus triggers estén activos, y el administrador no puede desactivarlos porque no es dueño de los objetos.
    \item \textbf{Carga masiva:} 60,000 registros (por encima del mínimo de 50,000) con un bloque PL/pgSQL por lotes con \texttt{COMMIT} intermedio.
    \item \textbf{Rendimiento:} tres consultas complejas analizadas con \texttt{EXPLAIN}; reducción del costo estimado de 95.4\,\% (Q1), 39.6\,\% (Q2) y 57.7\,\% (Q3), y del tiempo real de entre 73\,\% y 83\,\%, con la recreación del índice de Q1 como índice de cobertura.
    \item \textbf{Advisor y parámetros:} vistas equivalentes a \texttt{DBA\_\allowbreak{}ADVISOR\_\allowbreak{}FINDINGS} y \texttt{DBA\_\allowbreak{}ADVISOR\_\allowbreak{}RECOMMENDATIONS} que detectaron 5,285,500 filas leídas secuencialmente y, tras indexar, cero lecturas secuenciales en las particiones con datos; ajuste de \texttt{work\_\allowbreak{}mem}, \texttt{maintenance\_\allowbreak{}work\_\allowbreak{}mem} y \texttt{random\_\allowbreak{}page\_\allowbreak{}cost}.
    \item \textbf{Almacenamiento:} monitoreo de tablespaces y particiones, creación de una partición nueva, reubicación del histórico a otro tablespace (3,020\,kB), procedimiento de mantenimiento programado para ejecutarse una vez al día con registro de cada corrida, y truncado de partición sin afectar el resto de la tabla.
    \item \textbf{Reproducibilidad:} todo el escenario se levanta desde cero con un archivo \texttt{docker-compose\allowbreak{}.yaml} y diez scripts, y las capturas de este documento se generan de la ejecución real.
\end{itemize}

\subsection{Actividades no cubiertas o con adaptación}

\begin{itemize}
    \item \textbf{\texttt{DB\_\allowbreak{}BLOCK\_\allowbreak{}SIZE}:} se consultó, pero PostgreSQL lo fija al compilar el motor (8192 bytes) y no puede modificarse; la afinación de memoria se hizo con \texttt{work\_\allowbreak{}mem}, \texttt{maintenance\_\allowbreak{}work\_\allowbreak{}mem} y \texttt{random\_\allowbreak{}page\_\allowbreak{}cost}. No se modificó \texttt{shared\_\allowbreak{}buffers} porque requiere reiniciar el servicio.
    \item \textbf{\texttt{AUTOEXTEND} y tamaño de datafiles:} no existen en PostgreSQL; el tamaño de cada tablespace se definió por cálculo en \texttt{capacidad\_\allowbreak{}tablespaces} y se vigila con una alerta al 80\,\%, pero el motor no impide escribir más allá (el límite físico sería el volumen), y el ajuste dinámico se hizo con particiones y con la reubicación de objetos entre tablespaces.
    \item \textbf{MAC nativo:} PostgreSQL no tiene Oracle Label Security; se implementó con niveles de autorización y \textit{Row Level Security}, que cubre el objetivo (control obligatorio por clasificación) pero con una implementación propia. El MAC aplica al analista y al auditor; el administrador y el superusuario quedan fuera por diseño.
    \item \textbf{Inmutabilidad de la bitácora:} los triggers bloquean cualquier modificación por SQL, pero un superusuario puede desactivarlos; la protección completa requiere copiar la bitácora a un servidor externo o usar \texttt{pgAudit}.
    \item \textbf{Advisor automático:} \texttt{DBA\_\allowbreak{}ADVISOR\_\allowbreak{}*} no existe; se reconstruyó con vistas sobre las estadísticas, por lo que las recomendaciones (crear/eliminar índices) las aplicó el equipo y no se ejecutan de forma automática; lo que sí corre solo es el mantenimiento de índices y estadísticas (sección~\ref{sec:mant_programado}).
    \item \textbf{Entorno:} se usó Docker en lugar de la máquina virtual VirtualBox mencionada en el enunciado, igual que en las prácticas 1 a 4; la presentación y las preguntas del profesor se pueden resolver ejecutando las consultas de la sección~\ref{sec:demo} sobre el mismo contenedor.
    \item \textbf{Alcance de la medición:} con 60,000 filas que caben en memoria, las diferencias en tiempo real son de milisegundos; el beneficio será mayor con volúmenes productivos, donde también pesará el costo de escritura de los índices.
\end{itemize}

\subsection{Reflexión}

El proyecto muestra que la seguridad, el rendimiento y el almacenamiento deben diseñarse juntos: la partición por fecha permitió podar lecturas, mover datos históricos a otro tablespace y truncar una partición completa; la separación en roles y niveles permitió dar a cada usuario solo lo necesario, y descubrimos que en PostgreSQL ser dueño de un objeto equivale a poder desactivar sus controles, por lo que el administrador no debe ser propietario; y las vistas del advisor convirtieron la afinación en un proceso medible (hallazgo, recomendación, aplicación y verificación) y no en prueba y error. También confirmamos que cada motor resuelve la misma necesidad de forma distinta y que conviene conocer sus límites (índices locales por partición, sin \texttt{AUTOEXTEND} ni \texttt{DB\_\allowbreak{}BLOCK\_\allowbreak{}SIZE} configurable) antes de migrar un diseño de Oracle a PostgreSQL. La propuesta para llevar este diseño a producción está en la sección~\ref{sec:recomendaciones}.

% ==========================================
% REFERENCIAS
% ==========================================
\section{Referencias}

\begin{thebibliography}{99}

\bibitem{carrasco2026proyecto}
Carrasco Zúñiga, G. G. (2026).
\emph{Proyecto Bases de Datos Avanzadas: Diseño, Seguridad y Optimización de un Sistema de Gestión de Empleados}, versión 1.
Facultad de Ingeniería, Universidad Nacional Autónoma de México (UNAM).

\bibitem{pg_partitioning}
The PostgreSQL Global Development Group. (2026).
\emph{PostgreSQL 16 Documentation: Table Partitioning}.
Recuperado de: \url{https://www.postgresql.org/docs/16/ddl-partitioning.html}

\bibitem{pg_rls}
The PostgreSQL Global Development Group. (2026).
\emph{PostgreSQL 16 Documentation: Row Security Policies}.
Recuperado de: \url{https://www.postgresql.org/docs/16/ddl-rowsecurity.html}

\bibitem{pg_tablespaces}
The PostgreSQL Global Development Group. (2026).
\emph{PostgreSQL 16 Documentation: Tablespaces}.
Recuperado de: \url{https://www.postgresql.org/docs/16/manage-ag-tablespaces.html}

\bibitem{pg_explain}
The PostgreSQL Global Development Group. (2026).
\emph{PostgreSQL 16 Documentation: Using EXPLAIN}.
Recuperado de: \url{https://www.postgresql.org/docs/16/using-explain.html}

\bibitem{pg_pgss}
The PostgreSQL Global Development Group. (2026).
\emph{PostgreSQL 16 Documentation: pg\_\allowbreak{}stat\_\allowbreak{}statements}.
Recuperado de: \url{https://www.postgresql.org/docs/16/pgstatstatements.html}

\bibitem{pg_config}
The PostgreSQL Global Development Group. (2026).
\emph{PostgreSQL 16 Documentation: Resource Consumption (work\_\allowbreak{}mem, shared\_\allowbreak{}buffers)}.
Recuperado de: \url{https://www.postgresql.org/docs/16/runtime-config-resource.html}

\bibitem{pg_rules_privileges}
The PostgreSQL Global Development Group. (2026).
\emph{PostgreSQL 16 Documentation: Rules and Privileges (vistas \texttt{security\_\allowbreak{}barrier})}.
Recuperado de: \url{https://www.postgresql.org/docs/16/rules-privileges.html}

\bibitem{pg_ddl_priv}
The PostgreSQL Global Development Group. (2026).
\emph{PostgreSQL 16 Documentation: Privileges (propiedad de objetos)}.
Recuperado de: \url{https://www.postgresql.org/docs/16/ddl-priv.html}

\bibitem{pg_reassign}
The PostgreSQL Global Development Group. (2026).
\emph{PostgreSQL 16 Documentation: REASSIGN OWNED}.
Recuperado de: \url{https://www.postgresql.org/docs/16/sql-reassign-owned.html}

\bibitem{pgaudit}
pgAudit Project. (2026).
\emph{pgAudit: PostgreSQL Audit Extension}.
Recuperado de: \url{https://www.pgaudit.org/}

\end{thebibliography}

\end{document}
